% update_ch1_review.tex — Драфт обновлений для part1.tex «Обзор литературы»
% На основе Introduction статьи INTENSE (01_introduction_ru.tex)
%
% ИНСТРУКЦИЯ: текст ниже содержит материал для двух подсекций Главы 1:
% 1) Подсекция 1.2 — «Современные представления о механизмах нейронного кодирования»
%    - 1.2.1 «Индивидуальный код» — place cells, mixed selectivity, experience-dependent tuning
%    - 1.2.3 «Смешанный подход» — проблема коррелированных переменных, идентифицируемость
% 2) Подсекция 1.3 — «Методы анализа многомерной нейронной активности»
%    - Новая подсекция: информационно-теоретические методы
%
% Уникальный контент (отсутствует в DRIADA):
% - Автоматизация поведения (MoSeq, DeepAction, MARS) → сотни признаков
% - Кинетика кальциевого индикатора → ложные корреляции
% - Деконволюция: модель-зависимые искажения
% - Ограничения линейных методов и декодеров для задачи селективности
%
% При внесении — проверить пересечение с DRIADA (part3.tex, секция 3.2):
% оба текста упоминают NIT, MINT, FRITES, CEBRA, dPCA.
% Решение: в Гл.1 — научная мотивация (почему нужна MI). В DRIADA — функциональное
% сравнение пакетов (таблица уже есть). Сократить повторы в Гл.3.2.

% ===================================================================
% БЛОК 1: Для секции 1.2.1 «Индивидуальный код»
% ===================================================================

% Нейроны кодируют информацию, избирательно представляя различные стимулы,
% поведенческие паттерны или когнитивные состояния посредством своей активности,
% формируя тем самым основу нейронного кода~\cite{DayanAbbott2001}. На популяционном
% уровне эти сигналы представляют собой многомерные переменные, определяющие
% поведение и когнитивные процессы~\cite{Averbeck2006}.

% Нейронная селективность динамически формируется под влиянием опыта. В гиппокампе
% клетки места способны быстро образовывать пространственные рецептивные
% поля~\cite{Sotskov2022}, причём в ряде случаев эти карты остаются стабильными на
% протяжении недель~\cite{Wilson1993, Thompson1990}. В то же время широко
% распространена смешанная селективность~\cite{TyeMiller2024}: один нейрон может
% отражать комбинации нескольких стимулов, поведенческих или когнитивных переменных.
% Нейроны префронтальной коры формируют смешанную селективность к задачно-релевантным
% переменным в ходе обучения~\cite{mix3, Mante2013}. Даже канонические примеры
% функционально специализированных нейронов могут кодировать множественные переменные:
% клетки места зачастую отражают не только местоположение, но и время, вознаграждение
% и траектории~\cite{Zheng2021, Lian2022}.

% ===================================================================
% БЛОК 2: Для секции 1.2.3 «Смешанный подход» — проблема идентифицируемости
% ===================================================================

% Проблема атрибуции становится особенно острой в натуралистических условиях,
% поскольку поведенческие переменные по своей природе коррелированы. Развитие методов
% автоматического анализа поведения — как неконтролируемых (MoSeq)~\cite{Weinreb2024},
% так и контролируемых (DeepAction~\cite{Harris2023deepaction},
% MARS~\cite{Segalin2021}) и основанных на правилах
% (BehaviorDEPOT)~\cite{Gabriel2022} — позволяет извлекать десятки и сотни
% поведенческих признаков из видеозаписей, порождая богатые, но сильно коррелированные
% наборы данных. Это создаёт проблему идентифицируемости: кажущаяся селективность к
% одной переменной может объясняться селективностью к коррелированному ковариату.
% Ненаправленные движения объясняют значительную долю активности по всей
% коре~\cite{Musall2019}, тета-модуляция затрудняет разделение кодирования скорости
% и фазы в гиппокампе~\cite{Vanderwolf1969, McNaughton1983}.

% ===================================================================
% БЛОК 3: Для секции 1.3 — новая подсекция: информационно-теоретические методы
% ===================================================================

% Кальциевый имиджинг дополнительно усложняет задачу анализа селективности. В отличие
% от спайков, кальциевый флуоресцентный сигнал является непрерывным и характеризуется
% сложной временнóй динамикой. Даже самые быстрые индикаторы (jGCaMP8) демонстрируют
% времена затухания от десятков до сотен миллисекунд~\cite{Zhang2023}. Кинетика
% индикатора приводит к временнóй свёртке, при которой флуоресцентный сигнал от
% более ранних событий ложно коррелирует с последующим поведением~\cite{Wei2020}.
% Многие методы анализа спайков предполагают точечно-процессные наблюдения и не
% переносятся непосредственно на темпорально фильтрованные флуоресцентные сигналы;
% деконволюция может вносить модельно-зависимые искажения~\cite{Rupprecht2021}.
%
% Линейные меры не способны обнаруживать нелинейные зависимости, снижение
% размерности скрывает вклад отдельных нейронов, а вероятностные декодеры могут
% предсказывать поведение, не позволяя при этом количественно определить, какие
% именно переменные объясняют активность данного нейрона~\cite{Kriegeskorte2019}.
%
% Информационно-теоретические подходы предоставляют принципиальный формализм для
% количественной оценки связей между нейронной активностью и поведением. Наиболее
% полный из существующих инструментов — Neuroscience Information Toolbox
% (NIT)~\cite{Maffulli2022} — предлагает множество оценок ВИ с коррекцией смещения,
% однако рассчитан на парадигмы с повторными пробами (trial-based). Непрерывные
% записи свободного поведения требуют сохранения временнóй автокорреляции, оценки
% оптимальных задержек и контроля ковариации поведенческих признаков. Другие пакеты
% ориентированы на дополняющие уровни анализа: MINT~\cite{Lorenz2025} — на передачу
% информации на популяционном уровне; FRITES~\cite{Combrisson2022} — на групповой вывод
% в электрофизиологических данных; IDTxl~\cite{IDTxl2019} — на сетевую связность;
% HOI~\cite{Neri2024} — на взаимодействия высших порядков. Метод GCMI~\cite{Ince2017}
% обеспечивает эффективное вычисление ВИ, однако не решает задач выбора задержки и
% разделения кодирования от корреляционных артефактов.

% ===================================================================
% ПРИМЕЧАНИЕ: весь текст выше — в виде LaTeX-комментариев, так как это материал
% для интеграции в существующие секции part1.tex. При внесении — раскомментировать
% нужные абзацы и вставить в соответствующие места.
% ===================================================================
